\section{部分分式与反常积分}
\label{sec:02-Calculus/04-Integration/04}

计算积分 \(\int\dfrac{dx}{x^2-1}\) 时，常规的换元法与分部积分法均不再适用。分子中缺少因子 \(2x\)，无法直接凑微分得到 \(\ln\abs{x^2-1}\)；若采用分部积分，反而会增加分子中 \(x\) 的次数。陷入困境的本质在于被积函数的代数结构：分母是两个一次因子的乘积，而我们熟知的是分母仅含单一因子的积分形式。

为此，需先重构其代数形式。利用因式分解 \(x^2-1=(x-1)(x+1)\)，尝试将其拆分为两个简单分式之和：

\[
\frac{1}{(x-1)(x+1)}=\frac{A}{x-1}+\frac{B}{x+1}.
\]

右边通分得分子为 \((A+B)x+(A-B)\)。对比等式两边系数，要求 \(A+B=0\) 且 \(A-B=1\)，解得 \(A=1/2\)、\(B=-1/2\)。由此可得：

\[
\int\frac{dx}{x^2-1}=\frac12\ln\abs{x-1}-\frac12\ln\abs{x+1}+C .
\]

\subsection{部分分式是一条化归路径}

上述方法称为\textbf{部分分式分解}（Partial Fraction Decomposition）。它并未引入新的积分理论，而是将复杂的有理函数化归为几种已知的基本积分形式：\(\int dx/(ax+b)\) 导出对数函数，\(\int dx/(x^2+a^2)\) 导出反正切函数，\(\int dx/(x-a)^k\) 导出幂函数。

其通用流程十分明确：当分子次数不低于分母次数时，先通过多项式除法将假分式化为一个多项式与一个真分式之和；接着将真分式的分母分解为一次因子与不可约二次因子的乘积；对于每一个单重一次因子设置项 \(A/(x-a)\)，对 \(k\) 重一次因子需设出 \(1\) 到 \(k\) 次幂的所有项；对每一个不可约二次因子设置项 \((Ax+B)/(x^2+bx+c)\)；最后通过通分并对比系数解出待定常数。例如，若分母包含 \((x-1)^2(x^2+1)\)，分解形式即为 \(A/(x-1)+B/(x-1)^2+(Cx+D)/(x^2+1)\)。后续计算均为规范的代数运算，求解路径清晰确定。

值得注意的有两点：第一，该方法依赖于分母能够进行因式分解，而高于四次的多项式未必能用根式表达其根，此时往往需要借助数值方法；第二，原函数中的对数项 \(\ln\abs{\cdot}\) 仅在不含奇点的连通区间内有效。例如 \(1/(x^2-1)\) 的原函数在 \((-1,1)\) 和 \((1,\infty)\) 上是相互独立的，其积分常数可分别独立选取。这一点在后续处理带有奇点与无界区间的积分时至观重要。

\subsection{将积分端点拓展至无穷}

定积分的经典定义要求积分区间是有限闭区间。在书写 \(\int_1^\infty dx/x^2\) 时，并非直接将 \(x=\infty\) 代入原函数（无穷大并不是可以参与运算的实数），而是采取极限思想：先选取有限上界 \(R\) 计算标准的定积分，再观察当 \(R \to \infty\) 时该积分值的趋向：

\[
\int_a^\infty f(x)\,dx\eqdef\lim_{R\to\infty}\int_a^R f(x)\,dx .
\]

若极限存在且有限，则称该反常积分\textbf{收敛}；反之则称其\textbf{发散}。该定义包含两个独立的前提：在每一个有限区间 \([a,R]\) 上的定积分必须存在，且极限过程必须收敛。这一定义严密地避免了对 ``\(\infty\)'' 进行非法代数运算的问题。

\subsection{收敛性取决于尾部衰减速度}

以基准积分族 \(\int_1^\infty dx/x^p\) 为例。当 \(p\ne1\) 时，\(\int_1^R x^{-p}\,dx=(R^{1-p}-1)/(1-p)\)；当 \(p=1\) 时，原函数为 \(\ln R\)。令 \(R\to\infty\) 可得：若 \(p>1\)，则 \(R^{1-p}\to0\)，积分收敛于 \(1/(p-1)\)；若 \(p=1\)，\(\ln R\to\infty\)；若 \(p<1\)，\(R^{1-p}\to\infty\)。因此，该积分族收敛的充要条件是 \(p>1\)。临界点恰好落在 \(1/x\) 上，而 \(1/x\) 本身属于发散的情形。

\begin{center}
\begin{tikzpicture}[x=1.15cm,y=2.3cm,font=\small,>=stealth]
  \fill[red!14] (1,0) -- plot[domain=1:6,samples=70] (\x,{1/\x}) -- (6,0) -- cycle;
  \fill[blue!28] (1,0) -- plot[domain=1:6,samples=70] (\x,{1/(\x*\x)}) -- (6,0) -- cycle;
  \draw[->] (0.6,0) -- (6.9,0) node[right] {\(x\)};
  \draw[->] (0.8,-0.06) -- (0.8,1.25);
  \draw[red!70!black,thick] plot[domain=1:6.6,samples=80] (\x,{1/\x});
  \draw[blue!70!black,thick] plot[domain=1:6.6,samples=80] (\x,{1/(\x*\x)});
  \node[red!65!black] at (3.7,0.44) {\(1/x\)};
  \draw[red!45] (3.7,0.36) -- (3.7,0.29);
  \node[blue!65!black] at (5.3,0.30) {\(1/x^2\)};
  \draw[blue!45] (5.3,0.22) -- (5.3,0.05);
  \draw[dashed,black!45] (1,0) -- (1,1.05);
  \node[below] at (1,-0.04) {\(1\)};
  \node[right,black!65,font=\footnotesize,align=left] at (6.5,0.62)
    {红色区域向右\\ 面积无界};
\end{tikzpicture}

\emph{两条曲线的函数值均趋于零，图形呈现 ``逐渐变薄'' 的特征。然而 \(1/x^2\) 下方的总面积收敛为 1，而 \(1/x\) 下方的面积随着 \(\ln R\) 无限增长。 ``被积函数趋于零'' 仅是积分收敛的必要条件，决定总面积是否有限的核心在于函数的\textbf{衰减速度}。}
\end{center}

对于有限区间上的奇点（瑕点），处理方式类似，只需将极限过程移至奇点附近：

\[
\int_0^1\frac{dx}{\sqrt x}
\eqdef\lim_{\varepsilon\to0^+}\int_\varepsilon^1 x^{-1/2}\,dx
=\lim_{\varepsilon\to0^+}\bigl(2-2\sqrt\varepsilon\bigr)=2 .
\]

一般地，反常积分 \(\int_0^1 x^{-p}\,dx\) 收敛当且仅当 \(p<1\)。此处的阈值条件与无穷区间的情形正好相反：无穷远处要求函数衰减足够快，而原点附近则要求函数发散速率足够慢。两者尺度颠倒，不等号方向自然相反。

若奇点位于积分区间内部，必须将区间切分为两段，并要求两个极限各自独立收敛，不允许两侧的 \(+\infty\) 与 \(-\infty\) 相互抵消。通过对称截断定义的\textbf{柯西主值}（Cauchy Principal Value）是另一种概念：当反常积分收敛时，主值与反常积分一致；但当反常积分发散而主值存在时，主值并不等同于通常意义上的反常积分。

\subsection{无法求出原函数时的判别方法}

在实际应用中，许多反常积分的被积函数无法表达为初等函数的原函数。例如 \(\int e^{-x^2}dx\) 不存在初等原函数，但反常积分 \(\int_0^\infty e^{-x^2}dx\) 确实收敛于一个有限实数。

此时可采用\textbf{比较判别法}：若在尾部满足 \(0\le f(x)\le g(x)\) 且 \(\int g(x)dx\) 收敛，则 \(\int f(x)dx\) 亦收敛；若满足 \(f(x)\ge g(x)\ge0\) 且 \(\int g(x)dx\) 发散，则 \(\int f(x)dx\) 亦发散。例如，当 \(x\ge1\) 时，\(0\le 1/(x^2+1)\le 1/x^2\)，由于右侧积分收敛，左侧积分必然收敛，整个过程无需求出 \(\arctan x\)。比较判别法仅关注尾部的数量级比较，不依赖有限区间上的细节；同时，它仅判定收敛与否，不给出具体的积分值。

\subsection{期望的存在性与临界阈值}

将上述临界条件引入概率论，会产生明确的实际意义。

函数 \(f(x)=1/x^2\)（\(x\ge1\)）在 \([1,\infty)\) 上的积分为 1，是一个合法的概率密度函数。从该分布中采样可以得到有限的实数值，但在计算其数学期望时，被积函数需乘上变量 \(x\)：

\[
\int_1^\infty x\cdot\frac{1}{x^2}\,dx=\int_1^\infty\frac{dx}{x},
\]

此时积分恰好落在发散的一侧。这表明：即使分布合法且每次样本值均有限，其理论期望仍可能不存在。乘上因子 \(x\) 将衰减指数从 2 降至 1，刚好跨过了收敛的临界线。

柯西分布（Cauchy Distribution）即是这一现象的经典范例。其密度函数按 \(1/x^2\) 量级衰减，虽然归一化积分为 1，但其均值积分发散。这一性质在实际数据中有直观体现：从柯西分布中抽样时，样本均值不会随着样本量增加而收敛，反而会被频繁出现的极端值拉动，导致大数定律失效。在收入分布、城市规模以及网络流量等重尾数据中也存在类似特征—— ``绝大多数样本数值较小'' 并不能保证期望存在，决定期望是否收敛的往往是分布的尾部特征。

\subsection{方法的边界与延伸}

部分分式分解提供了有理函数积分的系统化解法，但其计算瓶颈在于多项式的因式分解；反常积分的比较判别法则提供了收敛性判断的有效手段，却不直接给出数值。参数 \(p=1\) 处于收敛与发散的临界线上，无论是 \(p\)-积分还是 \(p\)-级数，它均属于发散侧。

值得注意的是，\(\int_1^\infty dx/x^p\) 的收敛条件与级数 \(\sum 1/n^p\) 的收敛条件完全一致。这并非巧合，两者本质上都在探讨 ``无穷多个趋于零的量累加后是否有限''。两者的判别法可以相互转化，在下一章\hyperref[sec:02-Calculus/05-Sequences-and-Series/03]{``正项级数的比较''}中，将对这一对应关系展开系统讨论。